\appendices
\section{Rule-Based CVE Parser for Syscall Extraction}
\label{appendix:cve_parser}

\subsection{Scope and Design Principles}
\label{subsec:appendix_scope}

We design a deterministic, rule-based parser (rather than ML-based NLP) to extract high-risk
syscalls from CVE descriptions and vendor advisories, ensuring:

\begin{itemize}
  \item \textbf{Transparency:} Fully auditable and explainable pattern matching.
  \item \textbf{Precision:} Uses a kernel-specific syscall lexicon to avoid false positives.
  \item \textbf{Context-awareness:} Filters based on exploit conditions, kernel versions,
        architecture, capabilities, and device constraints.
\end{itemize}

The rule-based approach is preferred over ML-based NLP because it provides higher determinism,
does not introduce false positives from model inference, and allows full inspection of each
filtering decision. This is particularly critical in security contexts where auditability and
explainability are paramount.

\subsection{Syscall Dictionary Construction}
\label{subsec:appendix_dictionary}

The syscall dictionary is built directly from kernel sources:

\begin{lstlisting}[float=htbp, language=bash, caption={Extracting syscall names from kernel sources.}]
# Generated from kernel v5.15 headers
grep -r 'SYSCALL_DEFINE' arch/x86/kernel/ \
    | awk -F'(' '{print $2}' \
    | cut -d, -f1 \
    | sort -u
# --> [accept, access, acct, add_key, ...]
\end{lstlisting}

Sources:
\begin{itemize}
  \item \texttt{/arch/x86/entry/syscalls/syscall\_64.tbl}
  \item \texttt{/include/linux/syscalls.h}
  \item Linux man-pages Section~II
\end{itemize}

This ensures that only valid, in-kernel syscalls are candidates for extraction and analysis. The
resulting dictionary contains the canonical names of 335~syscalls supported by Linux kernel~6.x,
together with common aliases (e.g., \texttt{getpid} and \texttt{sys\_getpid}).

\subsection{Extraction and Filtering Rules}
\label{subsec:appendix_extraction_rules}

Regex patterns are applied to CVE descriptions, vendor advisories, and security mailing list
disclosures to identify syscall mentions and filter based on context.

\textbf{Primary regex patterns:}

\begin{lstlisting}[float=htbp, language=Python, caption={Regex patterns for identifying syscall mentions in CVE descriptions.}]
# Pattern 1: Direct syscall mention (e.g., "via the ptrace() syscall")
SYSCALL_DIRECT = re.compile(
    r'\b(?:the\s+)?'
    r'(?P<syscall>' + '|'.join(SYSCALL_DICT) + r')'
    r'(?:\(\d*\))?\s*(?:system\s+call|syscall)',
    re.IGNORECASE
)

# Pattern 2: Function name in exploit context
SYSCALL_EXPLOIT = re.compile(
    r'(?:exploit|abuse|invoke|call|use)\s+'
    r'(?:the\s+)?'
    r'(?P<syscall>' + '|'.join(SYSCALL_DICT) + r')'
    r'(?:\(\d*\))?',
    re.IGNORECASE
)

# Pattern 3: C annotation pattern (e.g., "setuid(0)" or "execve('/bin/sh', ...)")
SYSCALL_CODE = re.compile(
    r'\b(?P<syscall>' + '|'.join(SYSCALL_DICT) + r')'
    r'\s*\([^)]*\)',
    re.IGNORECASE
)
\end{lstlisting}

\textbf{Contextual filtering rules:}

After extracting candidate syscalls, the parser filters them based on:
\begin{enumerate}
  \item \textit{Kernel version applicability:} Version ranges are extracted from patterns like
        ``kernel versions prior to 5.12'' and compared against the host kernel version.
  \item \textit{Architecture:} Mentions like ``x86\_64 only'' or ``ARM-specific'' are matched and
        non-matching architectures are excluded.
  \item \textit{Capability requirements:} CVEs requiring \texttt{CAP\_SYS\_ADMIN} or other
        elevated capabilities are excluded if the container lacks that capability.
  \item \textit{Device constraints:} Exploits requiring specific hardware devices (e.g.,
        \texttt{/dev/kvm}, specific USB devices) are excluded if the device is unavailable.
\end{enumerate}

\subsection{Algorithm Pseudocode}
\label{subsec:appendix_pseudocode}

\begin{lstlisting}[float=htbp, language=Python, caption={Rule-based CVE parser: syscall extraction algorithm.}]
import re

def extract_high_risk_syscalls(cves, host_kernel, host_arch):
    """
    Extract the set of high-risk syscalls from CVE descriptions.

    Parameters:
        cves: List of CVE objects with a .desc field
        host_kernel: Host kernel version (tuple, e.g., (6, 8, 0))
        host_arch: Host architecture (str, e.g., "x86_64")

    Returns:
        S_blocked: Set of syscall names to block
    """
    S_blocked = set()

    for cve in cves:
        # Filter 1: Check kernel version applicability
        if not kernel_version_applicable(cve.desc, host_kernel):
            continue

        # Filter 2: Skip if architecture does not match
        if architecture_mismatch(cve.desc, host_arch):
            continue

        # Extract syscall mentions via regex
        found_syscalls = re.findall(syscall_regex, cve.desc)

        for syscall in found_syscalls:
            # Filter 3: Skip if special capability required
            #           but container does not have it
            if requires_special_cap(cve.desc) and \
               not container_has_cap():
                continue

            # Filter 4: Skip if specific device required
            #           but not available
            if needs_specific_device(cve.desc) and \
               not device_available():
                continue

            # Syscall passed all filters -> add to blocked set
            S_blocked.add(syscall)

    return S_blocked
\end{lstlisting}

\subsection{Real-World Filtering Examples}
\label{subsec:appendix_examples}

\textbf{Example~1: CVE-2022-0185~\cite{cve_2022_0185} (Linux Kernel --- fsconfig Overflow)}
\begin{itemize}
  \item \textit{CVE Description:} ``A heap-based buffer overflow was found in the Linux kernel's
        legacy\_parse\_param function in the Filesystem Context functionality\ldots allows a
        local user with CAP\_SYS\_ADMIN to escalate privileges\ldots''
  \item \textit{Syscalls Extracted:} \texttt{ioctl}, \texttt{mount}
  \item \textit{Filter Decision:} Requires \texttt{CAP\_SYS\_ADMIN} $\Rightarrow$ retained in
        $S_{\text{blocked}}$ if container has the capability; discarded otherwise.
  \item \textit{Outcome:} For unprivileged containers, this CVE is filtered out. For containers
        with \texttt{CAP\_SYS\_ADMIN}, \texttt{mount} is added to $S_{\text{blocked}}$.
\end{itemize}

\textbf{Example~2: CVE-2021-22555~\cite{cve_2021_22555} (NetFilter --- heap overflow)}
\begin{itemize}
  \item \textit{CVE Description:} ``A heap out-of-bounds write affecting Linux since
        v2.6.19-rc1 was discovered in net/netfilter/x\_tables.c\ldots can be exploited with
        setsockopt() to achieve local privilege escalation.''
  \item \textit{Syscall Extracted:} \texttt{setsockopt}
  \item \textit{Version Condition:} Affects kernels $< 5.12$; discarded if host kernel is
        5.15+.
  \item \textit{Filter Decision:} If host kernel is 5.15, this CVE is skipped and
        \texttt{setsockopt} is \textit{not} added to $S_{\text{blocked}}$, avoiding unnecessary
        blocking of legitimate network functionality.
\end{itemize}

These two examples illustrate how the context-aware CVE parser avoids both false positives
(blocking safe syscalls in certain environments) and false negatives (missing dangerous syscalls
when applicable). This deterministic approach ensures that $S_{\text{blocked}}$ accurately
reflects the true CVE attack surface of each deployed container environment.
